\section{Canonical convenient orientations and small failure-source sets}
\label{sec:small-failure-sources}

\noindent\emph{Draft status: this argument has undergone internal
adversarial review; it has not undergone external peer review.}

In this section $D$ is an arbitrary finite, nonempty oriented graph.
No strong-connectivity, minimality, or residual hypothesis is imposed.
All neighbourhoods and directed distances in the definitions below
refer to the original $D$. Put
\[
 U_D(x)=V(D)\setminus\bigl(\{x\}\cup N_D^+(x)\cup N_D^{++}(x)\bigr),
 \qquad g(x)=d_D^+(x)-d_D^{++}(x).
\]
In particular, a target in $U_D(x)$ is neither a first nor an exact
second neighbour of $x$.

Following Ghazal~\cite[Definition~1]{Ghazal2012}, a missing pair $ab$
is good if at least one of its orientations is convenient. We write
$G(a\to b)$ for the statement
\[
 \forall x\in V(D)\setminus\{a,b\}:\quad
 x\to a\Longrightarrow b\in N_D^+(x)\cup N_D^{++}(x).
\]
For any missing pair define its directional failure sets by
\[
 \begin{aligned}
 F_{ab}&=\{x\in V(D)\setminus\{a,b\}:x\to a,\ b\in U_D(x)\},\\
 F_{ba}&=\{x\in V(D)\setminus\{a,b\}:x\to b,\ a\in U_D(x)\}.
 \end{aligned}
\]
It is bad precisely when both sets are nonempty.

Define the protection relation by
\[
 y\triangleleft p
 \quad\Longleftrightarrow\quad
 y\to p\ \text{ and }
 N_D^-(y)\subseteq N_D^-(p)\setminus\{y\}.
\]
It implies $y\in U_D(p)$: a path $p\to z\to y$ would require
$z\to p$, contradicting $p\to z$. Let
\[
 \mathcal F(x)=\{y\in U_D(x):y\not\triangleleft x\},
 \qquad t(x)=|\mathcal F(x)|.
\]
The relation $\triangleleft$ is a strict partial order. Indeed,
$y\triangleleft p$ gives $d_D^-(y)<d_D^-(p)$, and if
$y\triangleleft p\triangleleft r$, then
$y\in N_D^-(p)\subseteq N_D^-(r)$ gives $y\to r$, and the
two containments compose to
$N_D^-(y)\subseteq N_D^-(r)\setminus\{y\}$.

Choose a linear order $\prec_0$ extending $\triangleleft$. Orient each
good missing pair by its unique convenient direction if only one
exists; if both exist, orient it from earlier to later in $\prec_0$.
We call these the canonical convenient orientations associated with
$\prec_0$.

\begin{lemma}[Compatibility of canonical convenient orientations]
\label{lem:canonical-convenient-compatibility}
Suppose $y\triangleleft p$ and both $pq$ and $qy$ are good missing
pairs. If the canonical orientation of $pq$ is $p\to q$, then the
canonical orientation of $qy$ is $y\to q$.
\end{lemma}
\begin{proof}
There are two convenient-direction implications:
\[
 G(p\to q)\Longrightarrow G(y\to q),
 \qquad
 G(q\to y)\Longrightarrow G(q\to p).
\]
For the first, every predecessor of $y$ is a predecessor of $p$.
For the second, every vertex reaching $y$ within two steps also
reaches $p$ within two steps: replace the last arc into $y$ by an arc
into $p$ using the protected in-neighbourhood containment. The relevant
sources are distinct from the missing-pair endpoints, as required
by convenience.

If the chosen directions were $p\to q$ and $q\to y$, these implications
would make both pairs convenient in both directions. The canonical
tie-breaking rule would force $p\prec_0 q\prec_0 y$, contrary to
$y\triangleleft p$ and hence $y\prec_0 p$.
\end{proof}

\begin{lemma}[At most three failure sources]
\label{lem:small-failure-sources}
Let $Z\subseteq V(D)$ satisfy $|Z|\le3$ and
\[
 F_{ab}\cup F_{ba}\subseteq Z
 \quad\text{for every bad missing pair }ab.
\]
Choose canonical convenient orientations using some linear extension
$\prec_0$ of protection. Let $K$ be the oriented graph on $Z$ whose
arcs are \emph{only} these added orientations of missing pairs inside
$Z$, and put
\[
 S=\{z\in Z:d_K^+(z)=0\}.
\]
If
\[
 t(z)\le g(z)+1\qquad(z\in Z\setminus S),
 \tag{SF1}\label{eq:small-source-budget}
\]
then $D$ has a Seymour vertex.
\end{lemma}

\begin{proof}
First, every missing pair inside $Z$ is good, so the definition of
$K$ accounts for every such pair. Otherwise its two endpoints and
two opposite failure sources would be four distinct vertices of $Z$.
Opposite sources are distinct, and are themselves non-adjacent: if
$u\in F_{ab}$ and $v\in F_{ba}$, either arc between them would create
one of the forbidden two-step paths.

No two vertices of $S$ form a missing pair, since its orientation in
$K$ would give one of them positive out-degree. For every bad pair,
at least one direction consequently has no failure source in $S$.
Indeed, if both directions had an $S$-source, those two opposite
sources would form a missing pair inside $S$. Orient each bad pair
in a direction avoiding failures from $S$, and complete the graph
to a tournament $T$ using these orientations and the canonical
orientations of all good pairs.

We record two path facts in the original graph. Every failure head
of a bad edge is unprotected: if $x\in F_{ab}$, an opposite witness
$v\in F_{ba}$ would, under $b\triangleleft x$, give $v\to x\to a$.
Thus
\[
 x\in F_{ab},\ ab\text{ bad}\quad\Longrightarrow\quad
 b\in\mathcal F(x).
 \tag{SF2}\label{eq:bad-failure-unprotected}
\]
Second, if the canonical convenient direction is $p\to q$ and
$y\triangleleft p$, then $G(p\to q)$ directly implies $G(y\to q)$.
Thus a missing pair $qy$ is automatically good, regardless of the
number of failure sources. Compatibility gives
\[
 y\triangleleft p,\quad qy\text{ missing}
 \quad\Longrightarrow\quad y\to q\text{ in }T.
 \tag{SF3}\label{eq:small-source-protected-direction}
\]

Choose an order $L$ maximizing the forward arcs of $T$, with feed
vertex $f$. When any $r$ added arcs leaving $f$ are reversed towards
$f$, all were backward in $L$. The forward-arc count of $L$ increases
by $r$, while that of every other order increases by at most $r$.
Thus $L$ remains a median order after either reversal used below.
The feed-vertex theorem~\cite{HavetThomasse2000} gives the SNP for
$f$ in the resulting tournament.

\emph{Case A: $f\notin Z\setminus S$.}
Reverse every added arc leaving $f$, obtaining $T'$. Then
$N_{T'}^+(f)=N_D^+(f)$. An exact two-step head is reached first
through an original out-neighbour. An original second arc or a
convenient added second arc reaches an original first or second
neighbour. A bad added second arc producing a genuinely new head
would be a failure from $f$. This is excluded if $f\notin Z$ by
the source-cover hypothesis, and if $f\in S$ by the bad-edge choices.
An added arc incident to $f$ cannot be the second arc of a two-step
path starting at $f$, because a tournament has no two-cycle.
Therefore
\[
 N_{T'}^{++}(f)\subseteq N_D^{++}(f).
\]
The SNP in $T'$ transfers to $D$, proving the conclusion.

\emph{Case B: $f\in Z\setminus S$.}
Choose an arc $f\to q$ of $K$. Reverse all other added arcs leaving
$f$, but retain this one. Then
\[
 N_{T'}^+(f)=N_D^+(f)\,\dot\cup\,\{q\}.
 \tag{SF4}\label{eq:small-source-retained-first}
\]
We claim the following exact-neighbourhood containment:
\[
 N_{T'}^{++}(f)
 \subseteq
 \bigl(N_D^{++}(f)\setminus\{q\}\bigr)
 \cup\bigl(\mathcal F(f)\setminus\{q\}\bigr).
 \tag{SF5}\label{eq:small-source-retained-envelope}
\]

Take an exact two-step head $y$ and a path $f\to a\to y$ in $T'$.
It is neither an original first neighbour nor $q$. If
$a\in N_D^+(f)$, an original or convenient second arc makes $y$
an original second neighbour. A bad added second arc with a new
head makes $y$ an unprotected original far target by
\eqref{eq:bad-failure-unprotected}. These added second arcs are not
incident to $f$, so their chosen orientations have not been reversed.

It remains to check the new intermediate $a=q$. The only issue is
a head $y\in U_D(f)$ protected from $f$. An original arc $q\to y$
would force $q\to f$ by protection, contrary to the original missing
pair $fq$. If $qy$ was missing, \eqref{eq:small-source-protected-direction}
forces its canonical direction to be $y\to q$. Neither endpoint
is $f$, so that arc is unchanged in $T'$. If $y\to q$ was original,
it too is unchanged. In all cases $q\to y$ is impossible. Thus every
new far head through $q$ is unprotected. Since $q$ is now first,
it must be removed from both parts of the envelope. This proves
\eqref{eq:small-source-retained-envelope}; it explicitly accounts
for all paths through the retained added first neighbour.

The original missing partner $q$ is either in $N_D^{++}(f)$ or in
$U_D(f)$. In the latter case it is unprotected, since it does not
point to $f$. Hence $q$ belongs to exactly one of the disjoint sets
$N_D^{++}(f)$ and $\mathcal F(f)$. Removing it in
\eqref{eq:small-source-retained-envelope} gives
\[
 d_{T'}^{++}(f)\le d_D^{++}(f)+t(f)-1,
 \qquad d_{T'}^+(f)=d_D^+(f)+1.
\]
The count is of distinct second-neighbourhood heads, not paths or
added arcs. By \eqref{eq:small-source-budget},
\[
 d_{T'}^+(f)-d_{T'}^{++}(f)
 \ge g(f)+2-t(f)\ge1,
\]
contradicting the feed-vertex theorem for the same median feed $f$.
Case B is impossible, while Case A gives the required Seymour vertex.
\end{proof}

The cardinality bound ensures that every internal missing pair of $Z$ is good. The protected-head compatibility itself has no source-cardinality restriction. All convenience and protection statements remain properties of the
original graph; no such property is transferred to an augmented graph
without the explicit path checks above.

\par\medskip\noindent\textbf{Relation to earlier completion arguments.}\quad 
Retaining one added out-arc at the feed vertex already appears in
Ghazal's proof of \cite[Theorem~3]{Ghazal2013}. That theorem
imposes a specific hypothesis on the entire missing graph. It does
not directly assert Lemma~\ref{lem:small-failure-sources} for a small
set of failure sources. The canonical-orientation compatibility and
the protected-head path check above are supplied explicitly here.
No comprehensive novelty or priority claim is made.
